\documentclass[a4paper,12pt]{article}

\usepackage[utf8]{inputenc}
\usepackage{graphicx}
\usepackage{hyperref}
\usepackage{listings}
\usepackage{amsmath}
\usepackage{amssymb}
\usepackage{geometry}
\geometry{margin=2.5cm}

\title{Relazione del Progetto \\ \large Titolo del Progetto}
\author{Nome Cognome}
\date{\today}

\begin{document}

\maketitle
\section{Scelta della modalità di comunicazione tra utenti e lavagna}

Per implementare la Kanban board condivisa, era necessario garantire una comunicazione affidabile e bidirezionale tra ciascun utente e la lavagna, che funge da peer centrale per la gestione dello stato delle schede. Abbiamo deciso di utilizzare i \textbf{socket TCP} per questa comunicazione, per i seguenti motivi:

\begin{itemize}
	\item \textbf{Affidabilità:} TCP garantisce che i dati inviati dall'utente alla lavagna e viceversa arrivino correttamente e nell'ordine previsto. Questo è essenziale per mantenere coerente lo stato della Kanban board tra tutti gli utenti.

	\item \textbf{Connessione persistente:} TCP permette di stabilire una connessione continua tra ciascun utente e la lavagna, evitando la necessità di ristabilire la connessione per ogni messaggio e consentendo aggiornamenti in tempo reale.

	\item \textbf{Semplicità di implementazione:} I socket TCP in C offrono un'interfaccia consolidata e ben documentata, facilitando lo sviluppo della comunicazione affidabile tra utenti e lavagna.

	\item \textbf{Gestione concorrente delle connessioni:} Con TCP è possibile gestire più utenti contemporaneamente tramite thread o meccanismi di multiplexing (\texttt{select()}, \texttt{poll()}), garantendo che tutti possano inviare e ricevere aggiornamenti in tempo reale.

	\item \textbf{Controllo del flusso e frammentazione automatica:} TCP gestisce internamente il controllo del flusso e la frammentazione dei messaggi, riducendo la complessità della gestione manuale dei dati tra utenti e lavagna.
\end{itemize}

In sintesi, la scelta dei socket TCP per la comunicazione utente–lavagna è stata motivata dalla necessità di garantire affidabilità, coerenza e aggiornamenti in tempo reale, elementi fondamentali per il corretto funzionamento della Kanban board condivisa.

\end{document}
